% Options for packages loaded elsewhere
% Options for packages loaded elsewhere
\PassOptionsToPackage{unicode}{hyperref}
\PassOptionsToPackage{hyphens}{url}
%
\documentclass[
  ignorenonframetext,
  aspectratio=169,
  russian,
]{beamer}
\newif\ifbibliography
\usepackage{pgfpages}
\setbeamertemplate{caption}[numbered]
\setbeamertemplate{caption label separator}{: }
\setbeamercolor{caption name}{fg=normal text.fg}
\beamertemplatenavigationsymbolshorizontal
% Prevent slide breaks in the middle of a paragraph
\widowpenalties 1 10000
\raggedbottom
\AtBeginPart{
  \frame{\partpage}
}
\AtBeginSection{
  \ifbibliography
  \else
    \frame{\sectionpage}
  \fi
}
\AtBeginSubsection{
  \frame{\subsectionpage}
}
\usepackage{iftex}
\ifPDFTeX
  \usepackage[T1]{fontenc}
  \usepackage[utf8]{inputenc}
  \usepackage{textcomp} % provide euro and other symbols
\else % if luatex or xetex
  \usepackage{unicode-math} % this also loads fontspec
  \defaultfontfeatures{Scale=MatchLowercase}
  \defaultfontfeatures[\rmfamily]{Ligatures=TeX,Scale=1}
\fi
\usepackage{lmodern}

\usefonttheme{serif} % use mainfont rather than sansfont for slide text
\ifPDFTeX\else
  % xetex/luatex font selection
  \setmainfont[]{DejaVu Serif}
  \setsansfont[]{DejaVu Sans}
  \setmonofont[]{DejaVu Sans Mono}
\fi
% Use upquote if available, for straight quotes in verbatim environments
\IfFileExists{upquote.sty}{\usepackage{upquote}}{}
\IfFileExists{microtype.sty}{% use microtype if available
  \usepackage[]{microtype}
  \UseMicrotypeSet[protrusion]{basicmath} % disable protrusion for tt fonts
}{}


\usepackage{longtable,booktabs,array}
\usepackage{calc} % for calculating minipage widths
\usepackage{caption}
% Make caption package work with longtable
\makeatletter
\def\fnum@table{\tablename~\thetable}
\makeatother
\usepackage{graphicx}
\makeatletter
\newsavebox\pandoc@box
\newcommand*\pandocbounded[1]{% scales image to fit in text height/width
  \sbox\pandoc@box{#1}%
  \Gscale@div\@tempa{\textheight}{\dimexpr\ht\pandoc@box+\dp\pandoc@box\relax}%
  \Gscale@div\@tempb{\linewidth}{\wd\pandoc@box}%
  \ifdim\@tempb\p@<\@tempa\p@\let\@tempa\@tempb\fi% select the smaller of both
  \ifdim\@tempa\p@<\p@\scalebox{\@tempa}{\usebox\pandoc@box}%
  \else\usebox{\pandoc@box}%
  \fi%
}
% Set default figure placement to htbp
\def\fps@figure{htbp}
\makeatother



\ifLuaTeX
\usepackage[bidi=basic,provide=*]{babel}
\else
\usepackage[bidi=default,provide=*]{babel}
\fi
\ifPDFTeX
\else
\babelfont{rm}[]{DejaVu Serif}
\fi
% get rid of language-specific shorthands (see #6817):
\let\LanguageShortHands\languageshorthands
\def\languageshorthands#1{}


\setlength{\emergencystretch}{3em} % prevent overfull lines

\providecommand{\tightlist}{%
  \setlength{\itemsep}{0pt}\setlength{\parskip}{0pt}}



 

\usepackage[]{csquotes}

\IfFileExists{plex-otf.sty}{
  %% Full TeXlive
  % \usepackage[%
  %   % math,
  %   RM={Scale=0.94},SS={Scale=0.94},SScon={Scale=0.94},TT={Scale=MatchLowercase,FakeStretch=0.9},DefaultFeatures={Ligatures=Common}
  % ]{plex-otf}
}{
  %% Local fallback fonts with Cyrillic support
  \usepackage{fontspec}
  \usefonttheme{professionalfonts}
  \AtBeginDocument{%
    \setmainfont{DejaVu Serif}
    \setsansfont{DejaVu Sans}
    \setmonofont{DejaVu Sans Mono}%
  }
}

%%% Load theme
% https://deic.uab.cat/~iblanes/beamer_gallery/
\IfFileExists{beamerthemegotham.sty}{
  %% Full TeXlive
  \usetheme{gotham}
  \gothamset{
    numbering=totalpagenumber,
    parttocframe default=off,
    sectiontocframe default=off,
    subsectiontocframe default=off,
  }
}{
  %% TinyTeX
  \usetheme{Madrid}
}
\makeatletter
\@ifpackageloaded{caption}{}{\usepackage{caption}}
\AtBeginDocument{%
\ifdefined\contentsname
  \renewcommand*\contentsname{Содержание}
\else
  \newcommand\contentsname{Содержание}
\fi
\ifdefined\listfigurename
  \renewcommand*\listfigurename{Список иллюстраций}
\else
  \newcommand\listfigurename{Список иллюстраций}
\fi
\ifdefined\listtablename
  \renewcommand*\listtablename{Список таблиц}
\else
  \newcommand\listtablename{Список таблиц}
\fi
\ifdefined\figurename
  \renewcommand*\figurename{Рисунок}
\else
  \newcommand\figurename{Рисунок}
\fi
\ifdefined\tablename
  \renewcommand*\tablename{Таблица}
\else
  \newcommand\tablename{Таблица}
\fi
}
\@ifpackageloaded{float}{}{\usepackage{float}}
\floatstyle{ruled}
\@ifundefined{c@chapter}{\newfloat{codelisting}{h}{lop}}{\newfloat{codelisting}{h}{lop}[chapter]}
\floatname{codelisting}{Список}
\newcommand*\listoflistings{\listof{codelisting}{Листинги}}
\makeatother
\makeatletter
\makeatother
\makeatletter
\@ifpackageloaded{caption}{}{\usepackage{caption}}
\@ifpackageloaded{subcaption}{}{\usepackage{subcaption}}
\makeatother

\usepackage{bookmark}
\IfFileExists{xurl.sty}{\usepackage{xurl}}{} % add URL line breaks if available
\urlstyle{same}
\hypersetup{
  pdftitle={Внешний курс №3},
  pdfauthor={Трусова Алина Александровна},
  pdflang={ru-RU},
  hidelinks,
  pdfcreator={LaTeX via pandoc}}


\title{Внешний курс №3}
\subtitle{Криптография на практике}
\author{Трусова Алина Александровна}
\date{}

\begin{document}
\frame{\titlepage}

\renewcommand*\contentsname{Содержание}
\begin{frame}[allowframebreaks]
  \frametitle{Содержание}
  \setcounter{tocdepth}{2}
  \tableofcontents
\end{frame}
\setcounter{tocdepth}{2}
\tableofcontents
}

\section{1. Цель
работы}\label{ux446ux435ux43bux44c-ux440ux430ux431ux43eux442ux44b}

\begin{frame}{1. Цель работы}
Пройти третий блок курса \enquote{Основы кибербезопасности}.
\end{frame}

\section{2. Выполнение блока 3: Криптография на
практике}\label{ux432ux44bux43fux43eux43bux43dux435ux43dux438ux435-ux431ux43bux43eux43aux430-3-ux43aux440ux438ux43fux442ux43eux433ux440ux430ux444ux438ux44f-ux43dux430-ux43fux440ux430ux43aux442ux438ux43aux435}

\begin{frame}{2.1 Введение в криптографию}
\phantomsection\label{ux432ux432ux435ux434ux435ux43dux438ux435-ux432-ux43aux440ux438ux43fux442ux43eux433ux440ux430ux444ux438ux44e}
\end{frame}

\begin{frame}
Для ответа на вопрос используется определение ассмиетричного шифрования
с двумя ключами (рис. {[}-@fig:001{]}).

\begin{figure}

{\centering \includegraphics[width=0.7\linewidth,height=\textheight,keepaspectratio]{../report/image/1.PNG}

}

\caption{Вопрос 4.1.1}

\end{figure}%
\end{frame}

\begin{frame}
Отмечены основные условия для криптографической хэш-функции (рис.
{[}-@fig:002{]}).

\begin{figure}

{\centering \includegraphics[width=0.7\linewidth,height=\textheight,keepaspectratio]{../report/image/2.PNG}

}

\caption{Вопрос 4.1.2}

\end{figure}%
\end{frame}

\begin{frame}
Отмечены алгоритмы цифровой подписи (рис. {[}-@fig:003{]}).

\begin{figure}

{\centering \includegraphics[width=0.7\linewidth,height=\textheight,keepaspectratio]{../report/image/3.PNG}

}

\caption{Вопрос 4.1.3}

\end{figure}%
\end{frame}

\begin{frame}
В информационной безопасности аутентификация сообщения или
аутентификация источника данных-это свойство, которое гарантирует, что
сообщение не было изменено во время передачи (целостность данных) и что
принимающая сторона может проверить источник сообщения (рис.
{[}-@fig:004{]})

\begin{figure}

{\centering \includegraphics[width=0.7\linewidth,height=\textheight,keepaspectratio]{../report/image/4.PNG}

}

\caption{Вопрос 4.1.4}

\end{figure}%
\end{frame}

\begin{frame}
Определение обмена ключами Диффи-Хэллмана. (рис. {[}-@fig:005{]}).

\begin{figure}

{\centering \includegraphics[width=0.7\linewidth,height=\textheight,keepaspectratio]{../report/image/5.PNG}

}

\caption{Вопрос 4.1.5}

\end{figure}%
\end{frame}

\begin{frame}{2.2 Цифровая подпись}
\phantomsection\label{ux446ux438ux444ux440ux43eux432ux430ux44f-ux43fux43eux434ux43fux438ux441ux44c}
\end{frame}

\begin{frame}
По определению цифровой подписи протокол ЭЦП относится к протоколам с
публичным ключом (рис. {[}-@fig:006{]}).

\begin{figure}

{\centering \includegraphics[width=0.7\linewidth,height=\textheight,keepaspectratio]{../report/image/6.PNG}

}

\caption{Вопрос 4.2.1}

\end{figure}%
\end{frame}

\begin{frame}
Алгоритм верификации электронной подписи состоит в следующем. На первом
этапе получатель сообщения строит собственный вариант хэш-функции
подписанного документа. На втором этапе происходит расшифровка
хэш-функции, содержащейся в сообщении с помощью открытого ключа
отправителя. На третьем этапе производится сравнение двух хэш- функций.
Их совпадение гарантирует одновременно подлинность содержимого документа
и его авторства (рис. {[}-@fig:007{]}).

\begin{figure}

{\centering \includegraphics[width=0.7\linewidth,height=\textheight,keepaspectratio]{../report/image/7.PNG}

}

\caption{Вопрос 4.2.2}

\end{figure}%
\end{frame}

\begin{frame}
Электронная подпись обеспечивает все указанное, кроме конфиденциальности
(рис. {[}-@fig:008{]}).

\begin{figure}

{\centering \includegraphics[width=0.7\linewidth,height=\textheight,keepaspectratio]{../report/image/8.PNG}

}

\caption{Вопрос 4.2.3}

\end{figure}%
\end{frame}

\begin{frame}
Для отправки налоговой отчетности в ФНС используется усиленная
квалифицированная электронная подпись (рис. {[}-@fig:009{]}).

\begin{figure}

{\centering \includegraphics[width=0.7\linewidth,height=\textheight,keepaspectratio]{../report/image/9.PNG}

}

\caption{Вопрос 4.2.4}

\end{figure}%
\end{frame}

\begin{frame}
Верный ответ укзаан на изображении (рис. {[}-@fig:010{]}).

\begin{figure}

{\centering \includegraphics[width=0.7\linewidth,height=\textheight,keepaspectratio]{../report/image/10.PNG}

}

\caption{Вопрос 4.2.5}

\end{figure}%
\end{frame}

\begin{frame}{2.3 Электронные платежи}
\phantomsection\label{ux44dux43bux435ux43aux442ux440ux43eux43dux43dux44bux435-ux43fux43bux430ux442ux435ux436ux438}
\end{frame}

\begin{frame}
Известные платежные системы - Visa, MasterCard, МИР (рис.
{[}-@fig:011{]}).

\begin{figure}

{\centering \includegraphics[width=0.7\linewidth,height=\textheight,keepaspectratio]{../report/image/11.PNG}

}

\caption{Вопрос 4.3.1}

\end{figure}%
\end{frame}

\begin{frame}
Верный ответ на изображении (рис. {[}-@fig:012{]}).

\begin{figure}

{\centering \includegraphics[width=0.7\linewidth,height=\textheight,keepaspectratio]{../report/image/12.PNG}

}

\caption{Вопрос 4.3.2}

\end{figure}%
\end{frame}

\begin{frame}
При онлайн платежах используется многофакторная аутентификация (рис.
{[}-@fig:013{]}).

\begin{figure}

{\centering \includegraphics[width=0.7\linewidth,height=\textheight,keepaspectratio]{../report/image/13.PNG}

}

\caption{Вопрос 4.3.3}

\end{figure}%
\end{frame}

\begin{frame}{2.4 Блокчейн}
\phantomsection\label{ux431ux43bux43eux43aux447ux435ux439ux43d}
\end{frame}

\begin{frame}
Proof-of-Work, или PoW, (доказательство выполнения работы) --- это
алгоритм достижения консенсуса в блокчейне; он используется для
подтверждения транзакций и создания новых блоков. С помощью PoW майнеры
конкурируют друг с другом за завершение транзакций в сети и за
вознаграждение. Пользователи сети отправляют друг другу цифровые токены,
после чего все транзакции собираются в блоки и записываются в
распределенный реестр, то есть в блокчейн. (рис. {[}-@fig:014{]}).

\begin{figure}

{\centering \includegraphics[width=0.7\linewidth,height=\textheight,keepaspectratio]{../report/image/14.PNG}

}

\caption{Вопрос 4.4.1}

\end{figure}%
\end{frame}

\begin{frame}
Консенсус блокчейна --- это процедура, в ходе которой участники сети
достигают согласия о текущем состоянии данных в сети. Благодаря этому
алгоритмы консенсуса устанавливают надежность и доверие к самоу сети.
(рис. {[}-@fig:015{]}).

\begin{figure}

{\centering \includegraphics[width=0.7\linewidth,height=\textheight,keepaspectratio]{../report/image/15.PNG}

}

\caption{Вопрос 4.4.2}

\end{figure}%
\end{frame}

\begin{frame}
Ответ - цифровая подпись (рис. {[}-@fig:016{]}).

\begin{figure}

{\centering \includegraphics[width=0.7\linewidth,height=\textheight,keepaspectratio]{../report/image/16.PNG}

}

\caption{Вопрос 4.4.3}

\end{figure}%
\end{frame}

\section{3. Выводы}\label{ux432ux44bux432ux43eux434ux44b}

\begin{frame}{3. Выводы}
Я прошла третий блок.
\end{frame}




\end{document}
